% Publication Materials for RAG Comparison Research
% Generated: January 10, 2026
% Note: Values marked with {{PLACEHOLDER}} will be filled after experiments complete

\documentclass{article}
\usepackage{booktabs}
\usepackage{amsmath}
\usepackage{graphicx}
\usepackage{hyperref}

\title{Comprehensive Evaluation of RAG Knowledge Bases for\\Smoking Cessation Counseling}
\author{Research Team}
\date{January 2026}

\begin{document}

\maketitle

\begin{abstract}
We present a rigorous comparison of Retrieval-Augmented Generation (RAG) approaches for smoking cessation counseling, using a novel independent test set validated for zero data leakage. Our analysis includes multiple metrics (ROUGE-L, BLEU, METEOR) across five independent trials with statistical significance testing.

\textbf{Key Finding:} [To be filled after experiments]

\textbf{Methodological Contribution:} We discovered severe data leakage in the original test set (7/15 questions with $\geq$90\% similarity to training data), invalidating previous results. This highlights the critical importance of test set independence verification in NLP research.
\end{abstract}

\section{Methodology}

\subsection{Data Leakage Discovery}
During our audit, we discovered that the original test set had significant overlap with training data:

\begin{table}[h]
\centering
\caption{Original Test Set Data Leakage Analysis}
\begin{tabular}{lcc}
\toprule
Severity Level & Count & Percentage \\
\midrule
Severe ($\geq$90\% similarity) & 7 & 47\% \\
Moderate (70--90\% similarity) & 6 & 40\% \\
Clean ($<$70\% similarity) & 2 & 13\% \\
\bottomrule
\end{tabular}
\label{tab:leakage}
\end{table}

Four test questions were \textit{exact matches} (100\% similarity) to questions in the training data, fundamentally invalidating the evaluation.

\subsection{Independent Test Set Creation}
We created a new 15-question test set with:
\begin{itemize}
    \item Maximum 50\% similarity to any training question (fuzzy string matching)
    \item Coverage of all original topics (benefits, cravings, NRT, relapse, etc.)
    \item Grammatically distinct formulations
    \item Domain expert-validated reference answers
\end{itemize}

\subsection{Evaluation Protocol}
\begin{itemize}
    \item \textbf{Model:} GPT-4o-mini (temperature=0.3, max\_tokens=200)
    \item \textbf{Retrieval:} ChromaDB with cosine similarity, k=3 documents
    \item \textbf{Runs:} 5 independent trials
    \item \textbf{Statistical Analysis:} 95\% confidence intervals, paired t-tests, Cohen's d
\end{itemize}

\section{Results}

\subsection{Main Results}

\begin{table}[h]
\centering
\caption{RAG Approach Comparison (5 runs, 15 questions each)}
\begin{tabular}{lcccc}
\toprule
Approach & ROUGE-L & BLEU & METEOR \\
\midrule
Baseline (No RAG) & {{BASELINE_ROUGEL}} & {{BASELINE_BLEU}} & {{BASELINE_METEOR}} \\
Raw Sources RAG & {{RAW_ROUGEL}} & {{RAW_BLEU}} & {{RAW_METEOR}} \\
AI-Generated RAG & {{AI_ROUGEL}} & {{AI_BLEU}} & {{AI_METEOR}} \\
Human-Curated RAG & {{HUMAN_ROUGEL}} & {{HUMAN_BLEU}} & {{HUMAN_METEOR}} \\
\bottomrule
\end{tabular}
\label{tab:main_results}
\end{table}

Values shown as mean $\pm$ 95\% CI.

\subsection{Statistical Comparisons (ROUGE-L)}

\begin{table}[h]
\centering
\caption{Pairwise Statistical Comparisons}
\begin{tabular}{lcccc}
\toprule
Comparison & Difference & Cohen's d & p-value & Significant \\
\midrule
AI vs Human & {{AI_VS_HUMAN_DIFF}}\% & {{AI_VS_HUMAN_D}} & {{AI_VS_HUMAN_P}} & {{AI_VS_HUMAN_SIG}} \\
AI vs Raw & {{AI_VS_RAW_DIFF}}\% & {{AI_VS_RAW_D}} & {{AI_VS_RAW_P}} & {{AI_VS_RAW_SIG}} \\
Human vs Raw & {{HUMAN_VS_RAW_DIFF}}\% & {{HUMAN_VS_RAW_D}} & {{HUMAN_VS_RAW_P}} & {{HUMAN_VS_RAW_SIG}} \\
\bottomrule
\end{tabular}
\label{tab:statistical}
\end{table}

\subsection{Professor's Targets Assessment}

\begin{table}[h]
\centering
\caption{Research Target Evaluation}
\begin{tabular}{lccc}
\toprule
Target & Expected & Actual & Status \\
\midrule
AI achieves 84\% of Human & $\geq$84\% & {{AI_HUMAN_RATIO}}\% & {{TARGET1_STATUS}} \\
AI 40\% better than Raw & $\geq$+40\% & {{AI_RAW_DIFF}}\% & {{TARGET2_STATUS}} \\
\bottomrule
\end{tabular}
\label{tab:targets}
\end{table}

\section{Discussion}

\subsection{Key Findings}
[To be filled after experiments]

\subsection{Why This Matters}
The discovery of data leakage in the original test set demonstrates:
\begin{enumerate}
    \item The importance of rigorous test set validation in NLP research
    \item That high performance scores may be artifacts of data contamination
    \item The need for fuzzy string matching during test set creation
\end{enumerate}

\subsection{Limitations}
\begin{itemize}
    \item Single domain (smoking cessation) - results may not generalize
    \item Single LLM (GPT-4o-mini) - other models may show different patterns
    \item 15 test questions - larger test sets would provide more statistical power
    \item Domain is well-covered in LLM training data - may reduce RAG benefit
\end{itemize}

\section{Conclusion}
[To be filled after experiments]

\section*{Acknowledgments}
We thank the professor for guidance and the research team for data curation efforts.

\end{document}
